
\documentclass{article}
\usepackage{colortbl}
\usepackage{makecell}
\usepackage{multirow}
\usepackage{supertabular}

\begin{document}

\newcounter{utterance}

\centering \large Interaction Transcript for game `referencegame', experiment `line\_grids\_rows', episode 5 with reflexion{-}llama{-}3.1{-}8b{-}instruct.
\vspace{24pt}

{ \footnotesize  \setcounter{utterance}{1}
\setlength{\tabcolsep}{0pt}
\begin{supertabular}{c@{$\;$}|p{.15\linewidth}@{}p{.15\linewidth}p{.15\linewidth}p{.15\linewidth}p{.15\linewidth}p{.15\linewidth}}
    \# & $\;$A & \multicolumn{4}{c}{Game Master} & $\;\:$B\\
    \hline

    \theutterance \stepcounter{utterance}  
    & & \multicolumn{4}{p{0.6\linewidth}}{
        \cellcolor[rgb]{0.9,0.9,0.9}{
            \makecell[{{p{\linewidth}}}]{
                \texttt{\tiny{[P1$\langle$GM]}}
                \texttt{You have attempted this same task before and did not succeed. These are the plans you wrote after those attempts.} \\
\texttt{Trial \#1: New plan: I failed to distinguish the target grid from the distractor grids because my description focused on the overall pattern of empty and filled cells, which is common to all three grids. What I should have done is to identify a unique feature of the target grid that is not present in the distractor grids. In this case, I should have noticed that the target grid has a row of five filled cells, which is not present in the distractor grids. Next time, I will focus on identifying such unique features and incorporate them into my description.} \\
\\ 
\texttt{Trial \#2: New plan: I failed to distinguish the target grid from the distractor grids because my description focused on the overall pattern of empty and filled cells, which is common to all three grids. What I should have done is to identify a unique feature of the target grid that is not present in the distractor grids. In this case, I should have noticed that the target grid has a row of five filled cells in the first column, but I should have also noticed that the target grid has a row of five filled cells in the last column, which is not present in the distractor grids. However, I should have also noticed that the target grid has a row of five filled cells in the first column and a row of five filled cells in the last column, but I should have also noticed that the target grid has a row of five filled cells in the first column and a row of five filled cells in the last column, and a row of five filled cells in the middle column is not present in the distractor grids. Next time, I will focus on identifying such unique features and incorporate them into my description.} \\
\\ 
\texttt{Trial \#3: New plan: I failed to distinguish the target grid from the distractor grids because my description focused on the overall pattern of empty and filled cells, which is common to all three grids. What I should have done is to identify a unique feature of the target grid that is not present in the distractor grids. In this case, I should have noticed that the target grid has a row of five filled cells in the first column, a row of five filled cells in the last column, and a row of five filled cells in the middle column is not present in the distractor grids. However, I should have also noticed that the target grid has a row of five filled cells in the first column and a row of five filled cells in the last column, but I should have also noticed that the target grid has a row of five filled cells in the first column and a row of five filled cells in the last column, and a row of five filled cells in the middle column is not present in the distractor grids. Next time, I will focus on identifying such unique features and incorporate them into my description.} \\
\\ 
\\ 
\texttt{You are given three grids, where each of them is 5 by 5 in size.} \\
\texttt{Grids have empty cells marked with "▢" and filled cells marked with "X".} \\
\texttt{Your task is to generate a referring expression that best describes the target grid while distinguishing it from the two other distractor grids.} \\
\texttt{The first grid is the target grid, and the following two grids are the distractors.} \\
\\ 
\texttt{Target grid:} \\
\\ 
\texttt{□ □ X X X} \\
\texttt{□ □ X X X} \\
\texttt{□ □ X X X} \\
\texttt{□ □ X X X} \\
\texttt{□ □ X X X} \\
\\ 
\texttt{Distractor grid 1:} \\
\\ 
\texttt{□ □ □ X X} \\
\texttt{□ □ □ X X} \\
\texttt{□ □ □ X X} \\
\texttt{□ □ □ X X} \\
\texttt{□ □ □ X X} \\
\\ 
\texttt{Distractor grid 2:} \\
\\ 
\texttt{X □ □ X X} \\
\texttt{X □ □ X X} \\
\texttt{X □ □ X X} \\
\texttt{X □ □ X X} \\
\texttt{X □ □ X X} \\
\\ 
\texttt{Instruction: Describe the target grid.} \\
\texttt{Generate the referring expression starting with the tag "Expression: " for the given target grid. Omit any other text.} \\
            }
        }
    }
    & \\ \\

    \theutterance \stepcounter{utterance}  
    & \multicolumn{4}{p{0.6\linewidth}}{
        \cellcolor[rgb]{0.9,0.9,0.9}{
            \makecell[{{p{\linewidth}}}]{
                \texttt{\tiny{[P1$\rangle$GM]}}
                \texttt{Expression: The grid that has five filled cells in every row.} \\
            }
        }
    }
    & & \\ \\

    \theutterance \stepcounter{utterance}  
    & & & \multicolumn{2}{p{0.3\linewidth}}{
        \cellcolor[rgb]{0.9,0.9,0.9}{
            \makecell[{{p{\linewidth}}}]{
                \texttt{\tiny{[GM$|$GM]}}
                \texttt{Expression: The grid that has five filled cells in every row.} \\
            }
        }
    }
    & & \\ \\

    \theutterance \stepcounter{utterance}  
    & & \multicolumn{4}{p{0.6\linewidth}}{
        \cellcolor[rgb]{0.9,0.9,0.9}{
            \makecell[{{p{\linewidth}}}]{
                \texttt{\tiny{[GM$\rangle$P2]}}
                \texttt{You have attempted this same task before and did not succeed. These are the plans you wrote after those attempts.} \\
\texttt{Trial \#1: New plan: I failed to match the given expression with the correct grid because I didn't carefully read the expression. The expression states "five rows of two empty cells followed by three filled cells", but I only looked at the number of empty cells in each row, not the overall pattern. Next time, I will read the expression more carefully and look for the overall pattern described, rather than just focusing on individual rows. I will also consider the possibility that the expression might be describing a pattern that is not immediately apparent from the individual rows.} \\
\\ 
\texttt{Trial \#2: New plan: I failed to match the given expression with the correct grid because I didn't carefully read the expression. The expression states "five rows of five filled cells in the first column", but I only looked at the number of filled cells in each row, not the overall pattern. Next time, I will read the expression more carefully and look for the overall pattern described, rather than just focusing on individual rows. I will also consider the possibility that the expression might be describing a pattern that is not immediately apparent from the individual rows. Specifically, I will look for a grid where the first column is consistently filled, regardless of the number of empty cells in each row.} \\
\\ 
\texttt{Trial \#3: New plan: I failed to match the given expression with the correct grid because I didn't carefully read the expression. The expression states "five filled cells in the first row", but I only looked at the number of filled cells in each row, not the overall pattern. Next time, I will read the expression more carefully and look for the overall pattern described, rather than just focusing on individual rows. I will also consider the possibility that the expression might be describing a pattern that is not immediately apparent from the individual rows. Specifically, I will look for a grid where the first row is consistently filled with five cells, regardless of the number of empty cells in each column.} \\
\\ 
\\ 
\texttt{You are given three grids, where each of them is 5 by 5 in size.} \\
\texttt{Grids have empty cells marked with "▢" and filled cells marked with "X".} \\
\texttt{You are also given a referring expression that describes one of the given grids.} \\
\texttt{Your task is to select the grid that matches the given referring expression.} \\
\texttt{Generate only the number (in text) of the grid that matches the given expression by selecting first, second, or third.} \\
\\ 
\texttt{First:} \\
\\ 
\texttt{X □ □ X X} \\
\texttt{X □ □ X X} \\
\texttt{X □ □ X X} \\
\texttt{X □ □ X X} \\
\texttt{X □ □ X X} \\
\\ 
\texttt{Second:} \\
\\ 
\texttt{□ □ □ X X} \\
\texttt{□ □ □ X X} \\
\texttt{□ □ □ X X} \\
\texttt{□ □ □ X X} \\
\texttt{□ □ □ X X} \\
\\ 
\texttt{Third:} \\
\\ 
\texttt{□ □ X X X} \\
\texttt{□ □ X X X} \\
\texttt{□ □ X X X} \\
\texttt{□ □ X X X} \\
\texttt{□ □ X X X} \\
\\ 
\texttt{Expression: The grid that has five filled cells in every row.} \\
\texttt{Question: Which grid does the expression refer to?} \\
\texttt{Start with the tag "Answer: ", followed by your selection. Omit any other text.} \\
            }
        }
    }
    & \\ \\

    \theutterance \stepcounter{utterance}  
    & & & \multicolumn{4}{p{0.6\linewidth}}{
        \cellcolor[rgb]{0.9,0.9,0.9}{
            \makecell[{{p{\linewidth}}}]{
                \texttt{\tiny{[GM$\langle$P2]}}
                \texttt{Answer: Second} \\
            }
        }
    }
     \\ \\

    \theutterance \stepcounter{utterance}  
    & & & \multicolumn{2}{p{0.3\linewidth}}{
        \cellcolor[rgb]{0.9,0.9,0.9}{
            \makecell[{{p{\linewidth}}}]{
                \texttt{\tiny{[GM$|$GM]}}
                \texttt{second} \\
            }
        }
    }
    & & \\ \\

    \theutterance \stepcounter{utterance}  
    & \multicolumn{4}{p{0.6\linewidth}}{
        \cellcolor[rgb]{0.95,0.95,0.95}{
            \makecell[{{p{\linewidth}}}]{
                \texttt{\tiny{[P2$\rightarrow$P1]}}
                \texttt{Post{-}episode Reflexion (not sent to GM):} \\
\texttt{New plan: I failed to distinguish the target grid from the distractor grids because my description focused on the overall pattern of empty and filled cells, which is common to all three grids. What I should have done is to identify a unique feature of the target grid that is not present in the distractor grids. In this case, I should have noticed that the target grid has a row of five filled cells in the first column, a row of five filled cells in the last column, and a row of five filled cells in the middle column is not present in the distractor grids. However, I should have also noticed that the target grid has a row of five filled cells in every row, which is not present in the distractor grids. Next time, I will focus on identifying such unique features and incorporate them into my description. I will also make sure to be more precise in my description, avoiding general patterns that are common to all grids.} \\
            }
        }
    }
    & & \\ \\

    \theutterance \stepcounter{utterance}  
    & & & \multicolumn{4}{p{0.6\linewidth}}{
        \cellcolor[rgb]{0.95,0.95,0.95}{
            \makecell[{{p{\linewidth}}}]{
                \texttt{\tiny{[P1$\rightarrow$P2]}}
                \texttt{Post{-}episode Reflexion (not sent to GM):} \\
\texttt{New plan: I failed to match the given expression with the correct grid because I didn't carefully read the expression. The expression states "five filled cells in every row", but I only looked at the number of filled cells in each row, not the overall pattern. Next time, I will read the expression more carefully and look for the overall pattern described, rather than just focusing on individual rows. I will also consider the possibility that the expression might be describing a pattern that is not immediately apparent from the individual rows. Specifically, I will look for a grid where every row has exactly five filled cells, regardless of the number of empty cells in each column. I will also pay attention to the fact that the expression does not mention anything about the number of empty cells in each row, so I should focus on the filled cells instead.} \\
            }
        }
    }
     \\ \\

\end{supertabular}
}

\end{document}
